\summarychapter{6}{The Rule of Books}{The Formation of the Bible}

\begin{inbrief}
Few decisions in Christian history were more basic than those of the first three centuries about the
Bible, the standard by which catholic Christianity stayed orthodox. The chapter explains how the church
kept the Old Testament, why it differs over the Apocrypha, and how twenty-seven Christian writings
came to be set apart as the New Testament.
\question{How did we get the Bible?}
\end{inbrief}

During the last great persecution, early in the fourth century, a Christian in Sicily was tried for
owning a copy of the Gospels. Ordered to read from them, he read the blessing on the persecuted and the
call to take up the cross, and was sent to his death. Roman officials had understood that to destroy
Christianity they had to burn its Scriptures. The Bible was the foundation of the church's evangelism,
teaching, worship and morality, the inspiration of its martyrs and its test against heresy: what do
the Scriptures say?

\ssection{The Basics of the Bible}
Jerome called the Scriptures ``the Divine Library,'' many books that are one; Greek speakers moved
from \emph{Biblia}, ``the Books,'' to ``the Book.'' The two Testaments are two covenants. A covenant is
a special relationship between two parties; in the Bible, the relationship between God and man that
God's grace begins and sustains. The old covenant was made with Abraham and with Israel, and the Old
Testament tells the story of the Jews and their worship; the new covenant was made by Jesus the
Messiah with his new people, and the New Testament tells of Christ and the birth of the church. The
Old promised, the New fulfilled. The church claimed the first along with the Jews and produced the
second in spite of them.

The word for the special standing of these books is \emph{canon}, Greek for a measuring rod: the list
of books the church discerned to be authoritative, inspired by God. Since the first Christians were
Jews, the church never lacked a canon. Jesus accepted the Old Testament as God's word, believed its
statements, obeyed its commands and set himself to fulfill it, and Christians inherited his attitude
and his ways of reading. Alexandrian Jews had read Scripture as \emph{allegory}, seeking a deeper
meaning beneath the literal; \emph{typology} finds a pattern linking two stories, so that Jesus is
read as a sacrifice or as a deliverer like Moses. By the third century scholars such as Origen of
Alexandria defended the Christian claim to the Old Testament by allegory, speaking of an obvious
meaning and a hidden one open only to those gifted by the Spirit. Pagan critics, whose case needed
the literal sense, were infuriated.

\ssection{The Question of the Apocrypha}
Keeping the Old Testament did not settle exactly which books it held. The Apocrypha, twelve or fifteen
books depending on how they are counted, are accepted by Roman Catholics and Eastern Orthodox and
rejected by most Protestants. The shorter list of thirty-nine books, arranged as Law, Prophets and
Writings and often called the ``Palestinian Canon,'' seems to be the one Jesus quoted. The Greek
translations made at Alexandria, chiefly the Septuagint, contained additional books and spread them
widely; this larger list is the ``Alexandrian Canon.'' Early Christians differed. In the West
Augustine accepted the Apocrypha, and the Roman Catholic Church followed him; in the sixteenth
century most Protestants rejected them.

From the start Christians had more than the old covenant: first the Word made flesh, then the living
apostles, whose message, spoken or written, was revered as the channel of the Lord's will. Churches
read letters from the apostles in worship and used accounts of the life of Jesus; the first gospels
were written no earlier than AD 60 or 70, but parts existed in writing before. How did twenty-seven
of these writings cross the line from important to holy? Shelley names three factors within the life
of the churches. First, the books had a self-evidencing quality, a power to transform lives: Justin
Martyr, after trying the Stoics, the Pythagoreans and the Platonists, was pointed by an old man on the
seashore to the prophets and became a Christian through reading Scripture, as did Tatian, Theophilus,
Hilary, Victorinus and Augustine. Second, books were used in worship. Paul asked that his letters be
exchanged and read, and Justin, in the first description of a Christian service, says that on the Day
of the Sun the ``memoirs of the apostles'' or the prophets were read and then explained. Yet reading
in church was not enough: the letter of Clement of Rome (about 96) was still read at Corinth eighty
years later and never entered the canon. Third, and most fundamental, was a tie to an apostle: was the
book written by an apostle or by someone in direct contact with their circle? As Clement put it,
the apostles come from Christ as Christ comes from God.

\ssection{A List of Christian Books}
Events forced the pace. About 140 Marcion, a rich shipowner from Sinope on the Black Sea and the son
of a bishop, came to Rome and fell under the influence of the gnostic Cerdo. He taught that the God of
the Old Testament was a God of wrath, the author of evil, who cared only for the Jews, while the God
of Jesus was grace and love for all. He therefore rejected the whole Old Testament and every Christian
writing that seemed to favor Jews or to contradict him, keeping only a cut-down Gospel of Luke and
ten letters of Paul, the one apostle who, in his view, had not corrupted the gospel; Paul, he said,
sits at Christ's right hand in heaven. Rome excommunicated him in 144, but Marcionite churches with
their own clergy and rites (they used no wine at communion) spread through Italy and as far as
Arabia, Armenia and Egypt, and villages near Damascus survived into the fourth century.

Marcion set the church two problems: a list of books shaped around Paul, and the loss of the Old
Testament. Paul had become, as Tertullian said, the apostle of the heretics, yet he meant too much to
be given up. The church restored the Pastoral Letters and the letters of the other apostles and
joined all the letters to four gospels, with Acts as the bridge. And it kept the Old Testament, for
two reasons. The God who creates is the God who redeems: the one who made the world chose Israel and
reclaims his creation through Jesus. And Christianity is a historical religion, resting on the belief
that at a particular time and place God took a hand in human affairs; faith therefore means facing
life's ``why, Lord?'' and still trusting God's good purpose. If Marcion pushed the churches toward
forming a New Testament, Montanus pushed them toward closing it.

\ssection{Fresh Voices from God}
Christianity is a religion of the Spirit; how then, in Tertullian's bitter phrase, was the Holy
Spirit ``chased into a book''? In the later second century enthusiasm was giving way to
ecclesiasticism. Prophecy was fading, numbers were growing, and the line between church and world was
blurring. Between 156 and 172 Montanus appeared in Asia Minor, demanding stricter discipline and
sharper separation from the world. Had he stopped there he would have done good. But he and his
prophetesses Prisca and Maximilla prophesied in ecstasy, as instruments on which the Paraclete
played, announced Christ's speedy return and called opposition blasphemy against the Spirit; churches
split. Montanus taught that a new age of the Spirit, with greater demands and higher revelations, had
surpassed the ages of the Law and of the Son. In the Spirit's name he thus denied that God's decisive
revelation had come in Jesus Christ.

To keep the gospel central the church set the apostolic writings apart as uniquely authoritative, the
standard for all later faith and practice. It did not stop believing in the Spirit: the Spirit who had
enabled the sacred books to be written now enabled people to understand, interpret and apply them.
The proof is the appearance of lists. The Muratorian Canon, written about 190 and published by
L.~A.~Muratori in 1740, names the four Gospels, Acts, thirteen letters of Paul, Jude, two letters of
John and the Revelation, together with two books that did not stay, the Wisdom of Solomon and the Apocalypse of Peter. By the early third century only a few books were questioned: Hebrews in the West, Revelation
in the East. Early in the fourth century Eusebius reported that some disputed James, Jude, 2 Peter and 2 and 3 John, and was puzzled by Revelation. The first list identical with today's is in the
Easter letter of Athanasius of Alexandria in 367; the councils of Hippo (393) and Carthage (397)
published the same list. (A chart in the chapter shows the New Testament as it gained acceptance.)

In one sense Christians created the canon; in another they only recognized the writings that had made
their authority felt. Their aim was to submit to the teaching of the apostles. The faith stayed
catholic, Shelley concludes, because it was apostolic.

\begin{sidebox}{The Logic of the Canon}
Everett Ferguson describes four steps in canonical thinking: believers recognize that some writings
carry divine authority; they begin to think of a boundary, since sacred writings differ from others
and may be limited in number; they move to a closed canon, an explicit list; and a consensus forms
around that list.
\end{sidebox}

\begin{keyfacts}
\keypeople{Jerome; Origen; Augustine; Justin Martyr; Clement of Rome; Marcion; Cerdo; Tertullian;
Montanus, Prisca and Maximilla; Eusebius; Athanasius; L.~A.~Muratori; Everett Ferguson.}
\keydates{%
  \kd{c. AD 60--70}{The first gospels written}
  \kd{c. 96}{Clement of Rome writes to Corinth}
  \kd{c. 140}{Marcion comes to Rome}
  \kd{144}{Marcion excommunicated}
  \kd{156--172}{Montanus begins to prophesy in Asia Minor}
  \kd{c. 190}{The Muratorian Canon lists the Christian books}
  \kd{367}{Easter letter of Athanasius: first list of the 27 New Testament books}
  \kd{393}{Council of Hippo lists the same 27 New Testament books}
  \kd{397}{Council of Carthage confirms the list of 27 books}}
\keyterms{Bible; covenant (testament); canon; allegory; typology; Apocrypha (deuterocanonical books);
Septuagint; Palestinian and Alexandrian canons; Marcionites; Montanism; Paraclete; Muratorian Canon.}
\end{keyfacts}
